% Part: incompleteness
% Chapter: theories-computability
% Section: q-is-ce

\documentclass[../../../include/open-logic-section]{subfiles}

\begin{document}

\olfileid{inc}{tcp}{qce}

\olsection{$\Th{Q}$ हा \printtoken{S}{c.e.}-संपूर्ण आहे}


\begin{thm}
$\Th{Q}$ हा !!{c.e.} आहे, पण निर्णेय नाही.
खरे तर तो संपूर्ण !!{c.e.} संच आहे.
\end{thm}

\begin{proof}
$\Th{Q}$ हा !!{c.e.} आहे हे पाहणे अवघड नाही:
तो अशा वाक्यांच्या संकेतांकांचा संच आहे की
$y$ हा वाक्याचा संकेतांक असेल, तर $\Th{Q}$ मध्ये
$y$ ने सांकेतिक केलेल्या वाक्याची एखादी सिद्धता $x$ अस्तित्वात असते:
\[
Q = \Setabs{y}{\lexists[x][\Prf[\Th{Q}](x,y)]}.
\]
पण $\Prf[\Th{Q}](x,y)$ संगणनक्षम आहे
(खरे तर आदिम पुनरावर्ती आहे), आणि वरच्या रूपात
लिहिता येणारा कोणताही संच संगणनक्षमपणे प्रगणनीय असतो.

तो संपूर्ण संगणनक्षमपणे प्रगणनीय संच आहे, असे म्हणणे
$K \leq_m Q$ असे म्हणण्यास समतुल्य आहे. येथे
$K = \Setabs{x}{\cfind{x}(x) \fdefined}$.
म्हणून $K$ चे~$\Th{Q}$ कडे न्यूनीकरण दाखवू.
क्लिनीचे विधेय $T(e,x,s)$ आदिम पुनरावर्ती असल्यामुळे
त्याचे $\Th{Q}$ मध्ये, समजा $!A_T$ या सूत्राने, निरूपण
होते. मग प्रत्येक $x$ साठी
\begin{align*}
x \in K & \lif \lexists[s][T(x,x,s)] \\
& \lif \lexists[s][(\Th{Q} \Proves !A_T(\num x,\num x, \num s))]
  \\
& \lif \Th{Q} \Proves \lexists[s][!A_T(\num x, \num x, s)].
\end{align*}
उलट $\Th{Q} \Proves \lexists[s][!A_T(\num x, \num x, s)]$
असेल, तर या सिद्धांताची स्वयंसिद्धके मानक नैसर्गिक-संख्या
प्रतिमानात सत्य असल्याने त्या प्रतिमानात हे अस्तित्ववाचक
वाक्य सत्य असते. म्हणून काही नैसर्गिक संख्या~$n$ साठी
$!A_T(\num x, \num x, \num n)$ सत्य असते.
$T(x,x,n)$ असत्य असेल, तर $!A_T$ हे $T$ चे
निरूपण करत असल्याने $\Th{Q}$ ही
$\lnot !A_T(\num x, \num x, \num n)$ सिद्ध करेल.
मात्र ते मानक प्रतिमानात असत्य सूत्र ठरेल आणि
$\Th{Q}$ ते सिद्ध करेल, हा विरोधाभास आहे.
म्हणून $T(x,x,n)$ सत्य असले पाहिजे; त्यावरून
$\cfind{x}(x) \fdefined$.

थोडक्यात, प्रत्येक $x$ साठी $x$ हा $K$ मध्ये
असतो तर आणि तरच $\Th{Q}$ हे
$\lexists[s][!A_T(\num x,\num x,s)]$ सिद्ध करते.
म्हणून $x$ पासून
$\lexists[s][!A_T(\num x,
  \num x, s)]$ या वाक्याचा संकेतांक देणारे
$f$ हे फलन $K$ चे~$\Th{Q}$ कडे न्यूनीकरण करते.
संख्यांकाचे प्रतिस्थापन आणि सूत्राचे सांकेतीकरण
संगणनक्षम असल्याने हे न्यूनीकरण संगणनक्षम आहे.
विकर्ण थांबण्याचा संच आधीच अनिर्णेय असल्याने
हे न्यूनीकरण या संचाचाही अनिर्णेयपणा दाखवते.
\end{proof}

\end{document}
